\section{Structural Inference}
\label{sec:test}

\subsection{Residual Law}

Define the residual statistic
\begin{equation}
    Q \;:=\; \sum_{j=1}^{L} \hat w_j \,\hat U_j^2,
    \label{eq:Q-stat}
\end{equation}
where $\hat w_j$ and $\hat U_j$ are produced by the two-step FWLS procedure in
\Cref{alg:fwls}.

\begin{theorem}[Exact Oracle Residual Law and Feasible Asymptotic Counterpart]
\label{thm:null-chi2}
Under $H_0$ and \Cref{ass:signal,ass:sigma,ass:lags}:
\begin{enumerate}[label=(\alph*)]
    \item \textup{(Oracle.)}  With true weights $w_j^*$, the statistic
    $Q^* = \sum_j w_j^* U_j^{*2}$, where $U_j^* = Y_j - \alpha - H x_j$
    are the residuals from the infeasible WLS projection,
    satisfies $Q^* \sim \chi^2(L-2)$ exactly.
    \item \textup{(Feasible.)}  The two-step statistic $Q$ satisfies
    $Q \dto \chi^2(L-2)$ as $n \to \infty$.
\end{enumerate}
\end{theorem}

\begin{proof}
\textbf{Part (a).}  Under $H_0$, $\mathbf{Y} = \mathbf{X}\boldsymbol\beta +
\mathbf{U}$ with $\mathbf{U} \sim \calN(\mathbf{0}, \sigma_0^2 n^{-1}
(\mathbf{W}^*)^{-1})$.  Define $\tilde{\mathbf{U}} =
(\mathbf{W}^*)^{1/2}\mathbf{U}/(\sigma_0 n^{-1/2})$, so $\tilde{\mathbf{U}}
\sim \calN(\mathbf{0}, \mathbf{I}_L)$.  The projection onto the column
space of $(\mathbf{W}^*)^{1/2}\mathbf{X}$ has rank 2, so the residual
projection has rank $L-2$.  Hence $Q^* = \norm{\tilde{\mathbf{U}} -
\mathbf{P}\tilde{\mathbf{U}}}^2 \sim \chi^2(L-2)$.

\textbf{Part (b).}  We have $Q = Q^* + (Q - Q^*)$.  By the linearisation
in the proof of \Cref{thm:clt-fwls}, $Q - Q^* = O_p(n^{-1/2})$,
so $Q \dto Q^*$ in distribution.  Since the $\chi^2$ CDF is continuous,
convergence in distribution passes through the CDF and $Q \dto \chi^2(L-2)$.
\end{proof}

\begin{remark}
    Part~(a) is exact because the oracle problem is a weighted Gaussian linear
    regression with known precision matrix.  Part~(b) is asymptotic because
    the feasible statistic uses estimated weights and the log-noise model is a
    delta-method approximation.  For finite $n$, the approximation error in
    \eqref{eq:log-noise} contributes a bias of order $O(n^{-1})$ per lag;
    this is negligible for the $\chi^2$ calibration at the rates we consider.
    Simulations confirm empirical sizes of 0.046--0.055 at $\alpha = 0.05$
    for all $L \in \{10,20,30,50\}$ and $n = 1000$ (see \Cref{tab:null}).
\end{remark}

\subsection{Testing Rule}

The residual rule at level $\alpha$ rejects $H_0$ when
\begin{equation}
    Q \;>\; q_{L-2,\,1-\alpha},
    \label{eq:test-rule}
\end{equation}
where $q_{L-2,\,1-\alpha}$ is the $(1-\alpha)$-quantile of $\chi^2(L-2)$.

For moderate $L$ (say $L \geq 20$), one may equivalently use the
standardised form
\[
    T \;:=\; \frac{Q - (L-2)}{\sqrt{2(L-2)}} \;\dto\; \calN(0,1),
\]
rejecting when $T > z_{1-\alpha}$.  However, the direct $\chi^2$ critical
value gives better finite-$L$ calibration (see \Cref{tab:null}).

\subsection{Behavior Under Departures}

\begin{proposition}[Power of the residual law]
\label{prop:alternative-response}
Under $H_1$ with parameter $\kappa > 0$, the marginal of each standardised
residual $\sqrt{w_j^*}\,\widehat U_j$ has variance $1 + \kappa^2/v_j$.
Consequently,
\[
    \E_\kappa[Q^*]
    = (L-2) + \kappa^2 \sum_{j=1}^L w_j^*
    = (L-2) + \frac{n\kappa^2}{\sigma_0^2}\sum_{j=1}^L j^{2H}.
\]
The oracle adequacy statistic is a variance-inflated quadratic form in normals,
not a non-central $\chi^2$ variable.  For local departures, the mean inflation
is governed by the same information scale as the lower bound.
\end{proposition}

\begin{proof}
The whitened residual vector has covariance $\mathbf I + \kappa^2 \mathbf W^*$.
Projecting onto the residual subspace of rank $L-2$ yields the stated mean
increment; the exact quadratic-form representation follows from diagonalising the
projected covariance. The distribution is therefore a weighted sum of centered
$\chi^2(1)$ terms.
\end{proof}

    \begin{remark}
    The alternative law of $Q^*$ is a variance-inflated quadratic form, not a
    non-central $\chi^2$.  Non-centrality arises from mean shifts, whereas
    \eqref{eq:Q-stat} under $H_1$ reflects covariance inflation.  The
    information-scale boundary and the oracle-equivalence result are unchanged.
    \end{remark}
